\documentclass[11pt]{article}
\usepackage[utf8]{inputenc}
\usepackage[T1]{fontenc}
\usepackage[brazilian]{babel}
\usepackage{amsmath, amssymb, amsthm}
\usepackage[margin=2.6cm]{geometry}
\usepackage{booktabs}
\newcommand{\code}[1]{\texttt{\small #1}}
\newcommand{\Z}{\mathbb Z}
\newcommand{\R}{\mathbb R}
\newcommand{\Q}{\mathbb Q}
\newcommand{\C}{\mathbb C}
\newcommand{\Hq}{\mathbb H}
\newcommand{\Oo}{\mathbb O}
\newcommand{\GF}{\mathrm{GF}}
\newcommand{\Zg}{\mathbb Z}
\usepackage{tcolorbox}
\newcommand{\medido}[1]{\smallskip\noindent\fbox{\parbox{0.97\linewidth}{\footnotesize
  \textbf{MEDIDO} --- evidência computacional, resíduo $0$: #1}}\smallskip}
\newcommand{\citado}[1]{\smallskip\noindent\fbox{\parbox{0.97\linewidth}{\footnotesize
  \textbf{CITADO} --- resultado da literatura, não demonstrado aqui: #1}}\smallskip}
\theoremstyle{plain}
\newtheorem{teorema}{Teorema}
\newtheorem{defn}{Definição}
\theoremstyle{remark}
\newtheorem{obs}{Observação}

\title{\textbf{A Multiplicação em $\R^n$}\\[3pt]
\large produto direto e produto cruzado, e onde fica a recursão}
\author{}\date{}

\begin{document}\maketitle

\begin{abstract}
\noindent A multiplicação do corpo autossimilar $\R^n$ escreve-se em \textbf{quatro peças}, e
delas só uma carrega a dimensão. As outras três --- a coordenada real, o produto interno e a parte
imaginária --- existem em toda dimensão e não mudam de forma. A quarta é o \textbf{produto
cruzado}, e é nela que fica a recursão: é ela que decide se o corpo comuta, se associa, e onde a
torre acaba. Os dois produtos --- \emph{direto} e \emph{cruzado} --- não são duas construções
avulsas: são as \textbf{duas metades} da decomposição de um bilinear qualquer em parte simétrica e
antissimétrica. E o mesmo par, uma escala acima, é a soma e a multiplicação do \emph{corpo de
corpos}. O que a dimensão $1,3,7$ limita é a \emph{norma multiplicativa}, e não a álgebra: no
espaço paramétrico $(p,r,t,s)$ os corpos clássicos são pontos especiais de um contínuo.
\emph{Quase} tudo o que aqui se afirma está medido, com o medidor indicado --- e o que
não está, por ser resultado conhecido da literatura, vai marcado como tal. A convenção que separa
os dois casos abre o texto, porque é ela que diz o peso de cada enunciado.
\end{abstract}

\section*{Convenção: o que é medido e o que é citado}

Este texto separa duas coisas que costumam aparecer misturadas, e a separação é deliberada:

\medido{o enunciado foi verificado por programa, com resíduo $0$ ou ao nível do erro de máquina,
e o medidor vai indicado. \emph{Isto é evidência computacional} --- prova que o enunciado vale
nos casos testados, e nada mais do que isso.}

\citado{o enunciado é um resultado conhecido, usado aqui e \emph{não} demonstrado. Vai dito onde
aparece, para que o leitor saiba de onde vem o peso.}

\noindent Onde não houver nenhuma das duas marcas, o enunciado é uma consequência algébrica
imediata do que ficou acima, e o leitor pode refazê-la em duas linhas.

\section{A construção}

Para cada $n$, um corpo:
\[
  \R^n \;=\; \GF(p^n) \;=\; \Z_p[x]/(p_n),\qquad p_n(x)=x^n-m\,x^{n-1}-1 ,
\]
com base $\{1,\sigma,\dots,\sigma^{n-1}\}$ --- as potências de \emph{um} gerador. Um elemento é uma
tupla, e a \textbf{notação algébrica} é a mesma tupla com o marcador dito:
\[
  x=(x_0,\dots,x_{n-1}) \;\longleftrightarrow\; x = x_0 + x_1\sigma + \dots + x_{n-1}\sigma^{n-1}.
\]
Não é abreviatura: é a passagem de um marcador \emph{posicional} para um marcador \emph{simbólico}.
E com ela as duas operações passam a ter a mesma cara --- ambas são as regras de polinómio --- e o
que sobra é a \textbf{borda}, a equação que baixa o grau:
\[
  \boxed{\;\sigma^n = m\,\sigma^{n-1} + 1\;}
\]
aplicada até o grau descer abaixo de $n$.

\begin{obs}[a borda é o parâmetro, e a família é a construção]
A lista $p_n(x)=x^n-mx^{n-1}-1$ é uma \emph{parametrização} --- a que dá os metais. A
\textbf{família} é $\Z_p[x]/(p)$ com $p$ irredutível, e trocar a borda dá outros corpos com a
\emph{mesma} recursão:
\begin{center}
\begin{tabular}{lll}
\toprule
$\sigma^2=m\sigma+1$ & os metais & $\Delta>0$, hiperbólico --- \emph{o gato} \\
$\sigma^2=-1$ & $\C$ & $\Delta=-4$, elíptico de ordem $4$ --- \emph{o esquilo} \\
$\sigma^2=\sigma-1$ & o sexto ciclotómico & $\Delta=-3$, elíptico de ordem $6$ \\
$\sigma^3=1$ & a raiz cúbica da unidade & ciclo de três \\
\bottomrule
\end{tabular}
\end{center}
Aplicando a fórmula recursiva a $n=2$ com $\sigma^2=-1$ sai \emph{exatamente} o produto complexo:
$(0,1)(0,1)=(-1,0)$, $(1,2)(1,-2)=(5,0)$, $(1,1)(1,1)=(0,2)$. \emph{Não é uma máquina parecida: é a
mesma fórmula, com o excedente a baixar por outra borda.} (\code{tools/algebra.c}, \code{hiper.c}.)
\end{obs}

\noindent Chama-se \textbf{família real} --- ou, o que é o mesmo, \textbf{cifra do rei} --- à base
ortonormal $\{1,\sigma,\dots,\sigma^{n-1}\}$: os $n$ marcadores, um por eixo.

\emph{Os nomes vêm do gerador.} Com $m=1$ a borda é $\sigma^2=\sigma+1$, cujo ponto fixo é
$\sigma=1+1/\sigma$ --- o número de ouro; a fração contínua dele é $[1;1,1,1,\dots]$, só termos
neutros, e é a essa sequência que se chama \textbf{cifra do rei}. As potências desse gerador
\emph{são} os eixos da base, e é a esse conjunto que se chama \textbf{família real}. Os dois nomes
olham o mesmo objeto de lados diferentes: a cifra é a \emph{sequência} que gera, a família é o
\emph{conjunto de eixos} que ela varre --- e é por isso que se usam como sinónimos.

\section{A multiplicação em quatro peças}

Separando o escalar do vetor, $z=(a_0,a)$ com $a\in\R^{n-1}$:
\[
 \boxed{\;(a_0,a)\cdot(b_0,b) \;=\; \bigl(\;
 \underbrace{a_0b_0}_{\text{coord.\ real}}-\underbrace{\langle a,b\rangle}_{\text{interno}}
 \;,\;\;
 \underbrace{a_0b+b_0a}_{\text{parte imaginária}}+\underbrace{a\times b}_{\textbf{cruzado}}
 \;\bigr)\;}
\]

\begin{teorema}[a fórmula dá $\C$ e $\Hq$ sem tabela de multiplicação]
Com o vetor de dimensão $1$ (cruzado inexistente) a fórmula reproduz o produto complexo; com
dimensão $3$ e o cruzado usual, reproduz a tabela de Hamilton --- $(1,2,3,4)(5,6,7,8) =
(-60,12,30,24)$. \emph{Não há tabela em lado nenhum: há a fórmula, e o que muda de um corpo para o
outro é a dimensão do vetor.}
\end{teorema}
\medido{\code{tools/rn.c} §R1 --- a fórmula aplicada a $\C$ (três pares) e a $\Hq$, conferida
contra a tabela de Hamilton.}

\noindent As quatro peças não têm o mesmo estatuto:

\begin{center}
\begin{tabular}{lll}
\toprule
\textbf{peça} & \textbf{existe em} & \textbf{o que faz} \\
\midrule
$a_0b_0$ & toda dimensão & a multiplicação de sempre \\
$\langle a,b\rangle$ & toda dimensão & simétrico; dá a norma \\
$a_0b+b_0a$ & toda dimensão & o escalar a agir no vetor \\
$a\times b$ & \textbf{só dim.\ $1$, $3$, $7$} & \textbf{carrega a dimensão} \\
\bottomrule
\end{tabular}
\end{center}

\noindent As três primeiras escrevem-se uma vez e servem sempre. \emph{A recursão fica na quarta.}

\section{Os dois produtos}

Não são listas de propriedades: são as duas metades de uma partição, e ela é única.

\begin{defn}[a partição]
Todo produto bilinear parte-se em duas, de um modo só:
\[
  B(a,b) \;=\; \underbrace{\tfrac12\bigl[B(a,b)+B(b,a)\bigr]}_{\text{simétrica: o \textbf{direto}}}
  \;+\; \underbrace{\tfrac12\bigl[B(a,b)-B(b,a)\bigr]}_{\text{antissimétrica: o \textbf{cruzado}}} .
\]
\end{defn}

\begin{center}
\begin{tabular}{ll}
\toprule
\textbf{PRODUTO DIRETO} --- a metade simétrica & \textbf{PRODUTO CRUZADO} --- a antissimétrica \\
\midrule
$\langle a,b\rangle=\langle b,a\rangle$: não vê a ordem & $a\times b=-b\times a$: \emph{é} a ordem \\
$\langle a,a\rangle\ge0$: dá a norma & $\langle a\times b,a\rangle=0$: perpendicular \\
devolve \textbf{escalar} --- baixa o grau & devolve \textbf{vetor} --- fica no espaço \\
existe em toda dimensão: sem obstrução & só em dim.\ $1,3,7$: \textbf{tem} obstrução \\
grupos: $G\times H$, os lados ignoram-se & grupos: $G\ltimes H$, um \textbf{age} no outro \\
$(g_1,h_1)(g_2,h_2)=(g_1g_2,\,h_1h_2)$ & $(g_1,h_1)(g_2,h_2)=(g_1g_2,\,g_1{\cdot}h_2+h_1)$ \\
\bottomrule
\end{tabular}
\end{center}

\begin{obs}[a diferença de fundo é uma só]
\emph{O direto não vê a ordem; o cruzado é a ordem.} As propriedades da tabela decorrem todas
desta distinção: o direto devolve escalar porque perdeu a direção --- que era justamente o que a
ordem guardava --- e o cruzado devolve vetor porque a guardou; o direto existe em toda dimensão
porque nada exige, e o cruzado tem obstrução porque exige um espaço onde caiba uma direção nova a
cada par. \emph{Os teoremas seguintes exploram exatamente essa separação.}
\end{obs}

\section{A recursão está no cruzado}

\begin{teorema}[a não-comutatividade \emph{é} o cruzado]
Das quatro peças, três não mudam ao trocar $a$ e $b$; só o cruzado muda, e muda de \emph{sinal}.
Logo
\[
  a\cdot b - b\cdot a \;=\; 2\,(a\times b),
\]
\emph{Tudo o que falta à comutatividade está no cruzado, e nada mais.}
\end{teorema}
\medido{\code{tools/rn.c} §R2 --- a diferença $ab-ba$ conferida contra $2(a\times b)$ nas quatro
coordenadas, e a simetria/antissimetria das peças verificada em separado.}

\begin{obs}[$\C$ comuta por causa da dimensão, não por escolha]
O vetor de $\C$ tem dimensão $1$, e o único antissimétrico em dimensão $1$ é o zero. \emph{O $i$
comuta consigo por não ter com quem cruzar.} A comutatividade é uma consequência da dimensão.
\end{obs}

\begin{teorema}[o que Hurwitz diz, e o que ele não diz]
Um produto bilinear antissimétrico que satisfaça $|a\times b|^2=|a|^2|b|^2-\langle a,b\rangle^2$
--- a identidade de Lagrange, que a \emph{norma multiplicativa} exige --- só existe em dimensão
$0$, $1$, $3$ e $7$. Daí:
\begin{center}
\begin{tabular}{llll}
\toprule
dim.\ do vetor & cruzado & corpo & \\
\midrule
$0$ & --- & $\R$ & comuta, associa \\
$1$ & é \emph{zero} & $\C$ & comuta, associa \\
$3$ & existe & $\Hq$ & \textbf{não} comuta; associa \\
$7$ & existe & $\Oo$ & não comuta \textbf{nem} associa \\
$15$ & \textbf{não há} & --- & divisores de zero: a torre \textbf{para} \\
\bottomrule
\end{tabular}
\end{center}
\end{teorema}
\citado{o teorema de Hurwitz (1898) sobre álgebras de composição. Usa-se e não se demonstra.}

\begin{obs}[e chamar-lhe \emph{obstrução} seria ler pela metade --- correção]
Uma versão anterior deste texto dizia ``a obstrução de Hurwitz''. Está errado, e o erro é de
leitura: \emph{Hurwitz não proíbe as outras álgebras --- classifica as que têm norma
multiplicativa.} São coisas diferentes. Uma álgebra em $\R^{16}$ existe e opera; o que ela não
tem é $\|zw\|=\|z\|\,\|w\|$, e a norma multiplicativa é uma \emph{exigência que se faz}, não uma
propriedade que o espaço deva ter.

E o próprio projeto já tinha a versão que mostra isso. Em \code{hiper/main.tex} a mesma fórmula
das quatro peças aparece \textbf{parametrizada}:
\[
  z\cdot_{p,r,t,s}w \;=\; \bigl(p\,z_0w_0 - r\langle a,b\rangle\bigr)
  \;+\; \bigl(t\,(z_0b+w_0a) + s\,(a\times b)\bigr),
\]
com um escalar por peça --- e o $s$ é exatamente o coeficiente do \emph{cruzado}. Nesse espaço,
$\R$, $\C$, $\Hq$ e $\Oo$ não são objetos discretos: são \textbf{pontos especiais de um contínuo}.
A norma multiplicativa vale precisamente onde $|p|=|r|=|t|=|s|=1$, e fora daí o produto continua a
existir --- só deixa de conservar a norma.
\medido{a norma multiplicativa verificada em $200$ pares por ponto: exata (erro $10^{-14}$) em
$(1,1,1,\pm1)$, e perdida assim que qualquer parâmetro sai de $\pm1$. E em \code{hiper/main.tex}
a órbita $\Z_2^3$ dá \emph{oito} representantes por classe, com erro $4{,}3\times10^{-14}$ ---
Hurwitz classifica as quatro \emph{classes de isomorfia}; a forma paramétrica mostra a estrutura
interna de cada uma.}
\emph{O que a dimensão $1,3,7$ limita é a norma, não a álgebra.} Dizer ``obstrução'' era eu a
tomar uma exigência minha por uma parede do mundo --- o mesmo erro que este projeto passa o tempo
a apanhar noutros sítios.
\end{obs}

\begin{obs}[e a hipótese que faz todo o trabalho é \emph{bilinear} --- os \emph{nne}-3D]
O enunciado acima diz ``um produto \textbf{bilinear} antissimétrico'', e está correto. O que eu
tirei dele --- \emph{que a torre para} --- omitia a hipótese, e omitir a hipótese transforma um
teorema numa proibição.

Gentil Lopes da Silva construiu em $\R^3$ uma multiplicação com \textbf{norma multiplicativa}, os
\emph{números não-negativos estendidos} (\emph{nne}-3D). Com $r_i=\sqrt{a_i^2+b_i^2}$ e
$\gamma=1-c_1c_2/(r_1r_2)$, e nos quatro casos conforme $r_1,r_2$ se anulem, o caso geral é
\[
  w_1\cdot w_2 = \bigl((a_1a_2-b_1b_2)\gamma,\;(a_1b_2+a_2b_1)\gamma,\;c_1r_2+c_2r_1\bigr).
\]
\medido{\code{tools/nne.c} --- os dois exemplos do livro reproduzidos exatos, e
$\|w_1w_2\|=\|w_1\|\,\|w_2\|$ com erro $10^{-15}$ em $500$ pares. \emph{Em $\R^3$.}}
E \emph{não há contradição nenhuma com Hurwitz}: esta multiplicação é homogénea mas
\textbf{não distributiva} --- usa a norma dos operandos, e a norma não é aditiva --- logo não é
bilinear, e está fora da hipótese do teorema.
\medido{\code{tools/nne.c} §N3 --- $w\cdot(u+v)$ difere de $w\cdot u+w\cdot v$ na terceira
coordenada; e $(\lambda w)\cdot u=\lambda(w\cdot u)$ exato.}
\begin{center}
\begin{tabular}{ll}
\toprule
\textbf{com} bilinearidade & $\R,\C,\Hq,\Oo$, e mais nada --- Hurwitz \\
\textbf{sem} bilinearidade & os \emph{nne}-3D, com norma multiplicativa em $\R^3$ \\
\bottomrule
\end{tabular}
\end{center}
\textbf{E generaliza a qualquer $n$ --- por recursão.} O livro diz que ``os $B$-3D generalizam, a
um só tempo, os números complexos e os \emph{nne}-2D'', com a imersão $(x,y)=(x,0,y)$; e olhando a
fórmula vê-se porquê: $(a_1a_2-b_1b_2,\;a_1b_2+a_2b_1)$ \emph{é} o produto complexo. Escrevendo
$z$ para o nível de baixo:
\[
  (z,c)\cdot(w,d) \;=\; \bigl(\,z\cdot w\cdot\gamma\;,\;\; c\|w\| + d\|z\|\,\bigr),
  \qquad \gamma = 1-\frac{cd}{\|z\|\,\|w\|},
\]
onde $z\cdot w$ é o produto do nível anterior. \emph{Sobe-se um nível de cada vez, e a norma
continua multiplicativa.}
\medido{\code{tools/nne.c} §N5 --- a recursão implementada e medida de $\R^2$ a $\R^7$: o erro vai
de $3{,}5\times10^{-15}$ a $1{,}4\times10^{-14}$, que é o arredondamento a acumular com os níveis.
E a recursão em $d=3$ dá o mesmo que a fórmula direta, conferido. \emph{Não há nível em que pare.}}

\textbf{E a interpretação --- que é o que faltava --- é que isto já estava escrito acima.} Comparando
com as quatro peças:
\begin{center}
\begin{tabular}{lll}
\toprule
\textbf{peça} & \textbf{bilinear} & \textbf{\emph{nne}} \\
\midrule
real & $a_0b_0-\langle a,b\rangle$ & $z\cdot w\cdot\gamma$ \\
imaginária & $a_0b+b_0a$ & $c\|w\|+d\|z\|$ \\
cruzado & $a\times b$, só em dim.\ $1,3,7$ & \emph{não há} --- o $\gamma$ ocupa o lugar \\
\bottomrule
\end{tabular}
\end{center}
\emph{O papel do escalar $a_0$ é feito pela norma $\|z\|$}, e dessa troca sai tudo o resto: a norma
não é linear, logo a multiplicação deixa de ser bilinear; e como não precisa do cruzado, não herda
a obstrução de dimensão. \textbf{O cruzado só existe em $1$, $3$ e $7$; a norma existe sempre.}
Trocar um pelo outro é trocar a bilinearidade pela liberdade de dimensão --- e é um preço, não um
almoço grátis: perde-se a distributividade, que é o que faz de um anel um anel.

\emph{As duas convivem: uma é dentro da hipótese, a outra é fora.} E é literalmente o erro que
este projeto apanha entrada a entrada no seu corpus --- a afirmação vale \emph{no corpo
declarado} --- cometido aqui sobre um teorema que eu próprio citara corretamente duas linhas
acima. \emph{(O material que consultei traz os \emph{nne}-2D e 3D explicitamente; o $7$D não vinha
escrito, mas sai da recursão acima e está medido.)}
\end{obs}
\medido{\code{tools/rn.c} §R4 --- a identidade de Lagrange verificada em dimensão $3$, que é a
condição que a norma multiplicativa impõe ao produto.}

\noindent Daqui decorre a razão de a recursão ficar no cruzado: \emph{subir de dimensão é
perguntar se ainda há para onde apontar}. O direto responde sempre afirmativamente, porque não
aponta; o cruzado é que tem de responder.

\begin{obs}[a recursão, a ordem e a indução são a mesma coisa, e moram no cruzado]
Vale dizer porquê, porque não é uma coincidência de três palavras. \emph{Induzir exige uma ordem}
--- o passo $n\to n+1$ só se dá onde há um antes e um depois. Das quatro peças, três não vêem a
ordem: $a_0b_0$, $\langle a,b\rangle$ e $a_0b+b_0a$ ficam iguais ao trocar $a$ por $b$. Só o
cruzado a vê, e é ele o único que muda de sinal. \emph{Logo a única peça capaz de carregar uma
indução é a que distingue o primeiro do segundo} --- e é a mesma que carrega a dimensão, a
não-comutatividade e o passo da torre. São quatro nomes para um lugar só.
\end{obs}

\section{O mesmo par, uma escala acima: o corpo de corpos}

Os elementos passam a ser corpos, e as duas operações são as mesmas duas.

\begin{teorema}[a soma direta não voa]
Em $\R^a\oplus\R^b$ os elementos $(1,0)$ e $(0,1)$ são não-nulos e o seu produto é $(0,0)$: há
divisor de zero, sem exceção. \emph{A dimensão fecha em $a+b$; a estrutura não fecha.} É o retrato
do produto direto --- põe dois na mesma gaiola, e nenhum age no outro.
\end{teorema}

\begin{teorema}[a lei que voa é o cruzado, balanceado sobre o comum]
\[
  \R^i\otimes_{\R^d}\R^j \;=\; \R^{\,\mathrm{lcm}(i,j)},\qquad d=\gcd(i,j),
\]
e o resultado é corpo, contém os dois pais, e é o \emph{menor} que os contém.
\end{teorema}

\begin{obs}[balancear sobre o comum \textbf{é} a ação]
A conta que separa as duas leituras é $a\cdot b=\mathrm{lcm}\cdot\gcd$: o tensorial \emph{cru} são
$\gcd$ cópias do balanceado. Sem balancear, cada lado conta o comum por si --- e contar duas vezes o
mesmo é precisamente \emph{não haver ação}. É por isso que o tensorial cru \emph{parece} ser a lei
quando só se testam espécies primas entre si: aí $\gcd=1$, há uma cópia só, e as duas coincidem.
Com divisor comum ele reparte-se, e o divisor de zero aparece.
\end{obs}
\medido{\code{tools/rn.c} §R6 e \code{tools/viveiro.c} --- o divisor de zero da soma direta por
varredura exaustiva; a identidade $a\cdot b=\mathrm{lcm}\cdot\gcd$ em todo par até $6$; e o filhote
a ter inverso para todo não-nulo.}

\begin{center}
\begin{tabular}{lll}
\toprule
$\oplus$ & o \textbf{direto} & simétrico, sem ação, soma as dimensões --- \emph{não voa} \\
$\otimes$ & o \textbf{cruzado} & com ação pelo comum, dá o $\mathrm{lcm}$ --- \emph{voa} \\
\bottomrule
\end{tabular}
\end{center}

\noindent \emph{``Não ver'' e ``agir'' são a mesma distinção nas duas escalas}: no vetor, o interno
não vê a ordem e o cruzado é a ordem; nos corpos, o direto não vê o outro e o cruzado age nele. E é
ela que separa somar de multiplicar.

\section{A régua, e a classificação}

A régua $(B,C)$ lê-se da própria borda: $\sigma^2=b_0+b_1\sigma$ é $\sigma^2-b_1\sigma-b_0=0$, logo
$B=-b_1$ e $C=-b_0$. \emph{Declarar o corpo é declarar a régua.}

\begin{teorema}[a classificação generaliza pela assinatura]
Em grau $n$ o invariante é a \textbf{assinatura} $(r,s)$ --- $r$ raízes reais, $s$ pares complexos,
$r+2s=n$ --- e há $\lfloor n/2\rfloor+1$ delas. Em grau $2$ são $(2,0)$ e $(0,1)$, mais o
degenerado da raiz dupla: \emph{exatamente} as três classes, e o $\Delta=B^2-4C$ é a assinatura
escrita num número --- o que só é possível porque ali ela só pode ser uma de duas. \emph{É o mesmo
corpo: nem sequer outra roupa, só outra interpretação.}
\end{teorema}
\medido{\code{tools/geral.c} §G5 --- a contagem de assinaturas em grau $2$ a $8$, e a assinatura de
seis polinómios conferida contra a esperada. As raízes reais contam-se por \emph{Sturm} --- a
sequência de restos de Euclides --- em aritmética exata, sem iteração.}

\begin{obs}[e o espectro misto não é uma classe nova]
$x^4-1$ tem duas raízes reais \emph{e} duas complexas --- mas ele \emph{fatoriza} em
$(x^2-1)(x^2+1)$: não é um corpo, são dois, um hiperbólico e um elíptico. É a mesma coisa do furo
em $n=5$, onde $x^5-x^4-1=(x^2-x+1)(x^3-x-1)$ separa o esquilo do gato. \emph{Um corpo não sobrevive
a ser dois.}
\end{obs}

\section{O que está medido}

\begin{center}
\begin{tabular}{ll}
\toprule
\code{tools/rn.c} & as quatro peças, os dois produtos, a obstrução, o corpo de corpos \\
\code{tools/algebra.c} & a álgebra global: o corpo é o argumento, e a borda o parâmetro \\
\code{tools/hiper.c} & a construção contra a leitura do $0/0$, e o furo em $n=5$ \\
\code{tools/geral.c} & grau $n$: a assinatura, Sturm, e $e^{At}$ contra a integração \\
\code{tools/zero.c} & as duas sementes da cifra, e o $J$ com $J^2=-I$ \\
\code{tools/viveiro.c} & o cruzamento, o \emph{lcm}, e o filhote a voar \\
\bottomrule
\end{tabular}
\end{center}

\noindent Todos com resíduo $0$ e integrados na bateria do repositório. O que é citado e não
medido --- o teorema de Hurwitz --- vai marcado onde aparece, e com a ressalva de que ele
\emph{classifica} as álgebras de norma multiplicativa, não proíbe as outras.

\appendix
\section{Os nomes, para quem abre este texto primeiro}

O projeto usa nomes próprios para três objetos que aparecem em toda parte. São apelidos de coisas
com nome técnico, e servem para não repetir a definição a cada uso:

\begin{center}
\begin{tabular}{lll}
\toprule
\textbf{apelido} & \textbf{o que é} & \textbf{como se reconhece} \\
\midrule
\textbf{gato} & o gerador que \emph{cresce} & $\Delta>0$, hiperbólico; autovalores reais, um deles \\
 & --- a \emph{companion} $\left[\begin{smallmatrix}m&1\\1&0\end{smallmatrix}\right]$
 & de módulo $>1$. Cresce e gasta; ordem infinita \\[3pt]
\textbf{esquilo} & o gerador que \emph{roda} & $\Delta<0$, elíptico; autovalores no círculo \\
 & --- uma rotação, $J=\left[\begin{smallmatrix}0&-1\\1&0\end{smallmatrix}\right]$
 & unitário. Gira e não gasta; ordem finita \\[3pt]
\textbf{rei} & o caso $m=1$ & a borda $\sigma^2=\sigma+1$; o número de ouro; \\
 & --- o número de ouro & a cifra $[1;1,1,1,\dots]$, só termos neutros \\
\bottomrule
\end{tabular}
\end{center}

\noindent E dois mais, que aparecem uma vez cada:

\begin{itemize}
\item \textbf{chicote}: o \emph{par} de raízes de uma borda quadrática. Elas somam o traço e
multiplicam o determinante, e por isso não se podem tratar separadamente --- ficar com uma é ficar
com metade da conta. No caso do rei, o par é $\varphi$ e $-1/\varphi$.
\item \textbf{borda}: a equação $\sigma^n=m\sigma^{n-1}+1$ que baixa o grau. É ela o parâmetro
que distingue um corpo do outro, e é ela que se declara quando se declara um corpo.
\end{itemize}

\noindent Nada neste texto depende dos apelidos: onde eles aparecem, o objeto técnico está dito ao
lado. Estão aqui porque o resto do projeto os usa, e porque nomear o que se repete poupa a
definição.

\end{document}
